\clearpage

\lehead[]{\normalfont\sffamily\hspace*{-2.00cm}\textcolor{white}{\colorbox{lightblue}{\makebox[1.60cm][r]{\thechapter}}}\hspace{0.17cm}\textcolor{lightblue}{\chaptertitle}}
\rohead[]{\textcolor{lightblue}{\chaptertitle}\normalfont\sffamily\hspace*{0.17cm}\textcolor{white}{\colorbox{lightblue}{\makebox[1.60cm][l]{\thechapter}}}\hspace{-2.00cm}}
%\chead[]{}
\rehead[]{\textcolor{lightblue}{AvHG, Inf, My}}
\lohead[]{\textcolor{lightblue}{AvHG, Inf, My}}

\lstset{style=myJava}


\section{Variablen}

Dass Computer Berechnungen wie

\begin{lstlisting}
2 + 5
12 * 7
\end{lstlisting}

ausführen können, wird euch nicht überraschen. Aber wie lassen sich solche
Ergebnisse für spätere Schritte merken? Und wie kann man eine Rechenvorschrift
so verallgemeinern, dass sie nicht nur für \lstinline|2 + 5|, sondern für beliebige
Zahlen funktioniert?

Dazu werden sogenannte \emph{Variablen} benutzt. Man kann sich eine Variable wie
eine beschriftete Schublade im \glqq Gedächtnis\grqq\ des Computers vorstellen, in
der ein Wert abgelegt und über ihren Namen jederzeit wieder abgerufen werden kann.

In sogenanntem \emph{Pseudocode} (einer gut lesbaren Mischung aus Alltagssprache
und Programmierstruktur) können wir damit etwa schreiben:

\begin{lstlisting}
a = 2
b = 5
summe = a + b
PRINT summe
\end{lstlisting}

\noindent \textbf{Wichtiger Unterschied zur Mathematik:} Das Gleichheitszeichen
\lstinline|=| bedeutet in der Informatik meist eine \textbf{Zuweisung} (von
rechts nach links): Der Wert auf der rechten Seite wird berechnet und in der
Variable auf der linken Seite gespeichert. Deshalb ist in der Programmierung auch
eine Zeile wie \lstinline|a = a + 1| völlig sinnvoll -- der alte Wert von
\lstinline|a| wird um $1$ erhöht und anschließend wieder unter dem Namen
\lstinline|a| abgelegt.


\section{Kontrollstrukturen}

Wenn Probleme so beschrieben werden sollen, dass sie als eindeutige
Handlungsanweisung Schritt für Schritt zu einem Ergebnis führen, reichen
einfache Anweisungen von oben nach unten (sogenannte \emph{Sequenzen}) oft nicht
aus. Man benötigt zusätzlich sogenannte \emph{Kontrollstrukturen}, die den
Ablauf eines Programms steuern.

Vermutlich habt ihr solche Kontrollstrukturen bereits unbewusst bei der Lösung
der ersten Aufgabe (der Quadratwurzel) benutzt: etwa um abhängig vom
vorangegangenen Ergebnis zu entscheiden, wie es weitergehen soll, oder um
bestimmte Rechenschritte so lange zu wiederholen, bis das Ergebnis genau genug
ist.

Obwohl es viele verschiedene Programmiersprachen gibt, verwenden fast alle
dieselben grundlegenden Kontrollstrukturen. Was ihr hier lernt, wird euch
deshalb auch dann von direktem Nutzen sein, wenn ihr später mit ganz anderen
Programmiersprachen arbeitet.


\subsection{Bedingte Anweisung und Verzweigung}

\emph{Bedingte Anweisungen} werden im Unterschied zu normalen Anweisungen nur
dann ausgeführt, wenn eine bestimmte Bedingung erfüllt (wahr) ist:

\begin{lstlisting}
FALLS (Ampelfarbe ist grün):
    Die Kreuzung überqueren.
\end{lstlisting}

Nur wenn die Bedingung zutrifft, wird der eingerückte Anweisungsblock
ausgeführt. Ist die Bedingung nicht erfüllt, wird der Block einfach
übersprungen.

Die bedingte Anweisung wird zur \emph{Verzweigung}, wenn es für den Fall, dass
die Bedingung nicht erfüllt ist, einen alternativen Anweisungsblock gibt:

\begin{lstlisting}
FALLS (Ampelfarbe ist grün):
    Die Kreuzung überqueren.
SONST:
    Anforderungsknopf drücken.
    An der Ampel warten.
\end{lstlisting}

In fast allen Programmiersprachen werden solche Verzweigungen mit den englischen
Schlüsselwörtern \lstinline|if| (falls) und \lstinline|else| (sonst) umgesetzt.


\subsection{Wiederholungsanweisungen (Schleifen)}

Als weitere zentrale Kontrollstruktur gibt es die \emph{Wiederholungsanweisung},
die meist kurz als \emph{Schleife} bezeichnet wird. Auch hier kontrolliert eine
Bedingung die Ausführung eines Anweisungsblocks. Im Unterschied zur bedingten
Anweisung wird der Block jedoch nicht maximal einmal ausgeführt, sondern so oft
wiederholt, wie die Bedingung erfüllt ist.

Mithilfe einer Schleife können wir die Handlungsanweisung für das Überqueren
einer Fußgängerampel realistisch vervollständigen:

\begin{lstlisting}
FALLS (Ampelfarbe ist grün):
    Die Kreuzung überqueren.
SONST:
    Anforderungsknopf drücken.
    WIEDERHOLE SOLANGE (Ampel nicht grün):
        An der Ampel warten.

    Die Kreuzung überqueren.
\end{lstlisting}
